\section{The MBPP without covariates}
\label{sec:no-covariates}

Here $\phi$ is a fixed convolution kernel. The MBPP equation~\eqref{eq:mbpp} is a
\emph{convolution} Volterra equation of the second kind, hence linear and
time-invariant (LTI), and everything is closed-form. This is the workhorse the
rest of the paper perturbs.

\subsection{Existence, uniqueness, and the resolvent}

Write the Volterra operator $(\mathcal V\xi)(t)=\int_0^t\phi(t-u)\xi(u)\dd u$, so
\eqref{eq:mbpp} reads $(\mathcal I-\mathcal V)\xi=s$.

\begin{theorem}[Well-posedness; see Appendix~\ref{app:volterra}]\label{thm:exist}
If $\phi$ is bounded on $[0,T]$ (true for the exponential kernel and for power-law
with positive offset) and $s\in C[0,T]$, then \eqref{eq:mbpp} has a unique
continuous solution on $[0,T]$, given by the Neumann series
$\xi=\sum_{n\ge0}\mathcal V^n s$.
\end{theorem}

\begin{intuition}
On a finite horizon, causal feedback can only compound so fast: the $n$-th echo
$\mathcal V^n$ carries a $t^n/n!$ factor that crushes it, so the loop
$\mathcal I-\mathcal V$ is \emph{always} invertible there, whatever the gain
$\kappa$. ``No runaway in finite time.'' Subcriticality ($\kappa<1$) is needed only
to extend this to $t=\infty$.
\end{intuition}

Separating the $n=0$ term, $\xi=s+h\conv s$ with the \emph{resolvent kernel}
$h=\sum_{n\ge1}\phi^{\conv n}$ ($\phi^{\conv n}$ the $n$-fold convolution). The
companion impulse response is $E=\delta+h$, so $\xi=E\conv s$
(Appendix~\ref{app:volterra}, \ref{app:transforms}).

\begin{proposition}[Global resolvent]\label{prop:resolvent}
If $\kappa=\int_0^\infty\phi<1$ then $h\in L^1(0,\infty)$ with
$\|h\|_1=\kappa/(1-\kappa)$, and $\xi=s+h\conv s$ holds on $[0,\infty)$.
\end{proposition}

\begin{intuition}
The resolvent sums the echo over \emph{all} offspring generations; it is globally
integrable exactly when each generation shrinks ($\kappa<1$), and its total mass
$\kappa/(1-\kappa)$ is the expected number of offspring per event. The resolvent
\emph{is} the renewal density of the branching tree.
\end{intuition}

\subsection{The exponential kernel closes everything}

For $\phi(t)=\kappa\theta e^{-\theta t}$ one finds (Appendix~\ref{app:ode})
\begin{equation}
h(t)=\sum_{n\ge1}\phi^{\conv n}(t)=\kappa\theta\,e^{(\kappa-1)\theta t},\qquad t>0,
\label{eq:h-exp}
\end{equation}
a single exponential (the series telescopes because one timescale stays one
timescale). The compensator obeys the \emph{same} equation with $s$ replaced by its
integral $\sfont S(t)=\int_0^t s$:
\begin{equation}
\Xi(t)=\sfont S(t)+\int_0^t\phi(t-u)\,\Xi(u)\dd u=\sfont S+h\conv\sfont S,
\qquad \Xi(x,y]=\Xi(y)-\Xi(x).
\end{equation}

\begin{example}[A single immigrant]
For $s=\delta(\cdot-a)$ (one immigrant at $a$), with $c:=(\kappa-1)\theta<0$,
\begin{equation*}
\xi(t)=\delta(t-a)+\kappa\theta\,e^{c(t-a)}\one[t>a],\qquad
\Xi^{\mathrm{endo}}(t)=\frac{\kappa}{1-\kappa}\bigl(1-e^{c(t-a)}\bigr)\one[t>a].
\end{equation*}
As $t\to\infty$, $\Xi^{\mathrm{endo}}\to\kappa/(1-\kappa)$: the mean offspring count
of one immigrant. The lone event leaves a \emph{decaying wake} of expected
offspring whose total mass is the branching sum.
\end{example}

\subsection{Multivariate MBPP}

With $M$ interacting components, $\xi\in\R^M$ solves the matrix convolution Volterra
equation
\begin{equation}
\xi(t)=s(t)+\int_0^t\Phi(t-u)\,\xi(u)\dd u,\qquad \Phi=(\phi_{mj}),
\label{eq:mbpp-mv}
\end{equation}
where $\phi_{mj}$ is how component $j$ excites component $m$. In the
transform domain (Appendix~\ref{app:transforms}) this is
$\hat\xi(\omega)=(I-\hat\Phi(\omega))^{-1}\hat s(\omega)$.

\begin{proposition}[Stationarity and spectral radius]\label{prop:stationary}
Under a constant baseline $s\equiv\bar\mu$ and stability $\rho(G)<1$ (where
$G_{mj}=\int\phi_{mj}=\hat\Phi(0)_{mj}$ is the branching matrix), the stationary
mean intensity is
\begin{equation}
\bar\xi=(I-G)^{-1}\bar\mu=\sum_{n\ge0}G^n\bar\mu .
\end{equation}
\end{proposition}
\begin{proof}
Stationarity in \eqref{eq:mbpp-mv} gives $\bar\xi=\bar\mu+G\bar\xi$; the Neumann
series converges iff $\rho(G)<1$.
\end{proof}

\begin{intuition}
The stationary intensity is the baseline amplified by the network's \emph{total}
feedback $(I-G)^{-1}=I+G+G^2+\cdots$: direct excitation, plus two-hop, plus
three-hop, \dots. The condition $\rho(G)<1$ is exactly what makes those multi-hop
echoes sum to something finite.
\end{intuition}

For exponential kernels the matrix equation is realised, with no transform, as a
finite \emph{linear ODE} in $M^2$ states $y_{mj}=\bigl(a_{mj}e^{-b_{mj}\cdot}\bigr)\conv\xi$
(Appendix~\ref{app:ode}):
\begin{equation}
y_{mj}'=a_{mj}\,\xi_j-b_{mj}\,y_{mj},\qquad \xi_m=s_m+\sum_j y_{mj},
\label{eq:statespace}
\end{equation}
with $a_{mj}=\kappa_{mj}\theta_{mj}$, $b_{mj}=\theta_{mj}$. This is the workhorse
solver \code{solve\_mbpp\_ode\_multivariate}; the numerical check
$\bar\xi=(I-G)^{-1}\bar\mu$ holds to machine precision.

\subsection{Fitting on interval-censored counts (IC-LL)}
\label{sub:icll}

Because the MBPP is Poisson (Theorem~\ref{thm:watanabe}), the bucket counts over an
observation grid $0=o_0<\cdots<o_m=T$ are \emph{independent}
$\mathrm{Poisson}(\Xi(o_{i-1},o_i])$. The negative log-likelihood is therefore a
simple sum.

\begin{proposition}[Interval-censored log-likelihood]\label{prop:icll}
For observed counts $\sfont C(o_{i-1},o_i]$,
\begin{equation}
\mathcal L_{\mathrm{IC}}(\theta)=\sum_{i=1}^m\Xi(o_{i-1},o_i;\theta]
-\sum_{i=1}^m \sfont C(o_{i-1},o_i]\,\log\Xi(o_{i-1},o_i;\theta],
\label{eq:icll}
\end{equation}
up to a constant in $\theta$. Minimising it over the kernel parameters recovers the
Hawkes parameters through the MBPP correspondence.
\end{proposition}

\begin{intuition}
Independent Poisson buckets $\Rightarrow$ the joint probability is a product
$\Rightarrow$ the log-likelihood is a sum of per-bucket Poisson terms in the
interval compensators $\Xi_i$. It compares a \emph{count} $\sfont C_i$ to an
\emph{expected count} $\Xi_i$ (both extensive), as it should. (Choosing a Gaussian
noise model instead of Poisson turns the loss into least squares; see
Appendix~\ref{app:estimation} for the Bregman/exponential-family picture and the
identifiability of $\kappa$ versus $\theta$.)
\end{intuition}

\begin{remark}[What is robustly estimable]
The branching ratio $\kappa$ (the gain) is well identified from counts; the
timescale $\theta$ is weakly identified, because a bucket integrates away the
within-bucket timing that carries $\theta$
(Appendix~\ref{app:estimation}). Treat $\kappa$ as trustworthy and report $\theta$
with wide intervals.
\end{remark}
